\documentclass[11pt]{article}
\usepackage[margin=27mm]{geometry}
\usepackage{amsmath,amssymb,amsthm,booktabs,hyperref}
\newtheorem{theorem}{Theorem}
\newtheorem{lemma}[theorem]{Lemma}
\newcommand{\Sz}{\operatorname{Sz}}
\newcommand{\Aut}{\operatorname{Aut}}
\title{Finite certificates for the rank-two geometries of $\Sz(8)$}
\author{Human Compression Project}
\date{}
\begin{document}
\maketitle
\begin{abstract}
The general classification reduces the remaining diagram calculation to four
families of subgroup pairs. We give certificates for the exhaustion and
isomorphism classes of these pairs, and for the incidence graphs they define.
The method uses the previously verified permutation model and subgroup
classes of $\Sz(8)$. It requires no new ambient enumeration or subgroup
classification. The independently verified literal certificates give $29$ outer classes and
$75$ inner classes. Two printed multiplicities are interchanged. We give the
exact association and compare all diagram entries with the source.
\end{abstract}

\section{Acquired data and the finite families}
Let $G=\Sz(8)$, of order $29120$. We acquire the complete permutation model
on $65$ points, its $22$ conjugacy classes of subgroups, their class lengths,
and a permutation $\phi$ inducing the field automorphism of order three.
Thus $\Aut(G)=G\rtimes\langle\phi\rangle$. These are premises supplied by the
previously reviewed subgroup treatment, rather than conclusions of the
calculation below. The general rank-two classification is acquired in the
same way.

Write a rank-two pair as $(M,L)$, and put $J=M\cap L$. The complete list of
families relevant to $G$ is
\[
\begin{array}{ccl}
 M&J&\text{possible }L\\\hline
 D_{14}&C_2&D_{14},D_{26},D_{10},C_2^2,C_4\\
 D_{14}&C_7&C_2^3\rtimes C_7\\
 C_{13}\rtimes C_4&C_4&C_{13}\rtimes C_4,C_5\rtimes C_4,C_4\times C_2\\
 C_5\rtimes C_4&C_4&C_5\rtimes C_4,C_4\times C_2.
\end{array}
\]
All three $G$-classes of $C_4\times C_2$ occur in the last two lines.
In each line $M$ is maximal in $G$. All eligible $J$ are conjugate in $M$:
this is Sylow conjugacy when $|J|=2$ or $4$, and uniqueness when $|J|=7$.
It therefore suffices to fix one embedded pair $M,J$ in each line.

\begin{lemma}\label{lem:incidence}
Suppose the subgroups of $G$ isomorphic to a fixed cyclic group $J$ form
one conjugacy class. Let $\mathcal H$ be one conjugacy class of subgroups,
of length $c(H)$, and let $z_J(H)$ be the number of subgroups of $H$
isomorphic to $J$. The number of members of $\mathcal H$ containing a
fixed $J$ is
\[
             \frac{c(H)z_J(H)}{c(J)}.
\]
\end{lemma}
\begin{proof}
Count the pairs $(K,J')$ with $K\in\mathcal H$ and $J'\le K$ conjugate to
$J$. Counting first by $K$ gives $c(H)z_J(H)$. Counting first by $J'$ gives
$c(J)$ times the desired number, since conjugation is transitive on the
class of $J$. Equating these counts proves the formula.
\end{proof}

The reader calculates $z_J(H)$ by listing the cyclic subgroups generated
by the elements of the acquired small group $H$ of order $|J|$. It checks
each proposed $L$ by a literal conjugating permutation, rejects repeated
embedded subgroups, and requires $M\cap L=J$. If $H$ has the same type as
$M$, the one forbidden subgroup $L=M$ is subtracted. Lemma~\ref{lem:incidence}
then proves exhaustion in each $G$-class separately. As $M$ is maximal
and $L\nleq M$, the groups $M,L$ generate $G$.

\section{Equivalence of the subgroup pairs}
Put $Q_G=N_G(M)\cap N_G(J)$. Since $G$ is simple and $M$ is proper maximal,
$N_G(M)=M$. Indeed, a larger normalizer would be $G$, making $M$ normal.
Consequently $Q_G=N_M(J)$ is found by inspecting at most $52$ elements.
Let $Q_A=N_{\Aut(G)}(M)\cap N_{\Aut(G)}(J)$.

\begin{lemma}\label{lem:outer}
If a supplied automorphism $a=g\phi$ normalizes both $M$ and $J$, then
$Q_A=\langle Q_G,a\rangle$.
\end{lemma}
\begin{proof}
The homomorphism $Q_A\longrightarrow\Aut(G)/G\cong C_3$ has kernel
$Q_G$. The image of $a$ generates $C_3$. Thus the subgroup generated by
$Q_G$ and $a$ has the same kernel and image as $Q_A$, and equals it.
\end{proof}

Conjugation by $Q_G$ determines equivalence of ordered pairs with fixed
$M,J$ under inner automorphisms. Adding conjugation by $a$ gives equivalence
under all automorphisms. Both actions are checked on the complete embedded
list of second groups.

When the parabolics have equal type, the two types may also be exchanged.
For each such $L$, the certificate supplies $t\in N_G(J)$ with $M^t=L$.
Such a $t$ exists: start with $M^g=L$, and conjugate the Sylow subgroup
$J^g$ to $J$ inside $L$. The exchanged pair $(L,M)$ is then sent by
$t^{-1}$ to $(M,M^{t^{-1}})$. We add this edge to the equivalence relation.
Every isomorphism either preserves or exchanges the two types, so these
edges and the normalizer actions generate the full relation. The number
of inner classes contained in each outer class is obtained from the two
finite component partitions. The outer quotient has order three, so this
number is one or three.

\section{The coset and graph certificate}
For each representative pair we use the right cosets $H\backslash G$,
where $H=M$ or $L$. Inversion identifies these with the source's left
cosets $gH$ and preserves incidence. The certificate gives a label $f_H(g)$
for every acquired element $g\in G$, and a permutation $A_s$ of the
vertex labels for every fixed ambient generator $s$.

\begin{lemma}\label{lem:cosets}
Suppose $f_H$ is surjective, $f_H^{-1}(f_H(1))=H$, and
\[
 f_H(gs)=A_s(f_H(g))\qquad(g\in G)
\]
for every ambient generator $s$. Then the fibres of $f_H$ are precisely
the right cosets of $H$.
\end{lemma}
\begin{proof}
Each $A_s$ is a permutation, so the displayed identity also holds for
inverse generators. It therefore extends to every word. A word equal to
$1$ acts trivially on each $f_H(g)$, and hence on every vertex by
surjectivity. We obtain a genuine $G$-action on the labels. It is transitive,
because each group element is a word in the generators, and its stabilizer
at $f_H(1)$ is exactly $H$ by the asserted identity fibre. The usual
orbit--stabilizer identification now gives the stated coset fibres.
\end{proof}

The incidence edges are exactly
\[
             \{(f_M(g),f_L(g)):g\in G\}.
\]
The reader requires their number to be $|G:J|$. Thus it constructs the
actual incidence graph, without relying on a coset library. Breadth-first
search from one vertex of each type gives the two point diameters, since
$G$ is transitive on each type. For the girth, remove the chosen vertex
of the first type and find the shortest path between any two of its
distinct neighbours; add two to the minimum length. Every shortest cycle
can be moved to pass through the chosen vertex. Conversely, each path
found in this manner closes to a cycle through it. This proves the
computed girth, rather than merely an upper bound.

The residue orders follow directly from the subgroup indices:
\[
(s_0,s_1)=([L:J]-1,[M:J]-1).
\]
Local double transitivity is checked on $J\backslash M$ and
$J\backslash L$. In a transitive action the point stabilizer is transitive
on all other points exactly when the action is doubly transitive. Here
that stabilizer is $J$, whose elements are already listed. These two
small orbit checks settle the local double-transitivity mark.

\section{The positive primitivity marks}
All pairs in the acquired family theorem are residually weakly primitive:
$J$ is maximal in both parabolics, and $M$ is maximal in $G$. It follows
that the geometry is residually primitive precisely when the second
parabolic is also maximal. This verifies the RPRI marks for the pairs
whose two parabolics belong to $D_{14}$, $C_{13}\rtimes C_4$ and
$C_5\rtimes C_4$ in the displayed families.

For the QPRI marks with second parabolic $D_{26}$ or $D_{10}$, the complete
acquired subgroup classification leaves only one possible proper overgroup
type, respectively $C_{13}\rtimes C_4$ or $C_5\rtimes C_4$. Such an
overgroup has a unique Sylow subgroup of order $13$ or $5$. It therefore
normalizes the rotation subgroup already contained in the dihedral group.
The acquired torus-normalizer theorem makes it the unique overgroup of
that type. In either case the subgroup interval is a single saturated
chain, since the remaining inclusion has index two.

The other positive QPRI case has second parabolic $H=C_2^3\rtimes C_7$.
Only a point stabilizer $B=P\rtimes C_7$ can be a proper maximal overgroup,
by the acquired list of maximal subgroups. The elementary abelian subgroup
of $H$ is the centre $Z$ of $P$. If another point stabilizer contained $H$,
its Sylow $2$-subgroup would contain a fixed nonidentity $z\in Z$.
All involutions in such a Sylow subgroup are central, and the acquired
centralizer formula $C_G(z)=P$ forces the same Sylow subgroup, hence the
same normalizer $B$. Finally $B/Z\cong C_2^3\rtimes C_7$, where the
complement acts irreducibly on $P/Z$. Subgroups between $H$ and $B$
correspond to invariant subspaces of $P/Z$, so there is no proper
intermediate subgroup. Again the interval from $H$ to $G$ is a single
saturated chain. These arguments verify the remaining positive QPRI
marks without a new subgroup search.

The marks on the $C_4\times C_2$ cases are corrected by the separately
accepted root group argument. An absent printed mark is not used as an
additional negative claim in the source comparison.

\section{Relation to the source}
The source compares pairs consisting of a geometry and its acting group
under automorphisms of $G$ [1, p.~100]. This is the equivalence just proved.
It is not an assertion that the underlying unmarked graphs are mutually
nonisomorphic. Repeated identical printed diagrams will be compared as
multisets of graph data and inner multiplicities. No artificial matching
between indistinguishable numbered occurrences is required.

The independently accepted corrections to diagrams 25 and 26 replace
the residue orders $(12,4)$ by $(4,12)$ in the printed orientation.
The accepted root group argument removes the QPRI marks from diagrams
27 and 29. These statements are kept separate from the graph calculation.
The printed graph data and multiplicities are not used by the producer.
The literal acquisition produces $29$ outer classes and $75$ inner classes.
After the already accepted corrections, $27$ printed entries match in full.
The remaining two entries require the following additional correction.

\begin{theorem}[Correction to the inner multiplicities]
In diagrams 25 and 26, the printed parenthetical multiplicities $3$ and $1$
are interchanged. Diagram 25, with edge triple $(6,4,6)$, has one inner class;
diagram 26, with edge triple $(5,3,6)$, has three.
\end{theorem}
\begin{proof}
Fix the acquired $M=C_{13}\rtimes C_4$ and its cyclic subgroup $J$ of order
four. Lemma~\ref{lem:incidence} gives exactly four second groups of type
$C_5\rtimes C_4$ containing $J$. In the certificate these have indices
$3,4,5,6$ (indices start at zero). The inner normalizer fixes each of them.
The supplied outer normalizer acts as the cycle $(3\ 5\ 6)$ and fixes $4$;
these are equalities of the explicit embedded subgroups. Thus the two
outer classes have inner multiplicities three and one respectively.
The certified coset graph for representative $3$ has point diameters
$(5,6)$ and girth $6$, giving edge triple $(5,3,6)$. The graph for
representative $4$ has point diameters $(6,6)$ and girth $8$, giving
$(6,4,6)$. This determines the association of the two multiplicities.
The native source page prints the opposite association.
\end{proof}

With this correction included, the complete multiset of the $29$
printed graph entries agrees with the computed multiset, including subgroup
labels, element counts, Borel orders, residue orders and local double
transitivity. The finite certificates and source transcription have distinct reviews.
The present application acquires those reviews; its proof and source
comparison are submitted together for independent mathematical review.

\begin{thebibliography}{9}
\bibitem{source} D. Leemans, \emph{The rank 2 geometries of the simple Suzuki
groups $\Sz(q)$}, Beitr\"age zur Algebra und Geometrie \textbf{39}
(1998), 97--120.
\end{thebibliography}
\end{document}
